\section{总结与延伸}

\subsection{全讲总结表}
\begin{table}[H]
\centering
\begin{tabular}{p{2.8cm}p{3.8cm}p{4.2cm}}
\toprule
\textbf{主题} & \textbf{核心观点} & \textbf{工程指向} \\
\midrule
概念框架 & Safety 与 Security 必须合并考虑 & 目标函数从输出安全升级为执行安全 \\
系统抽象 & Agent 多维能力带来组合攻击面 & 设计期就要做跨维度威胁建模 \\
攻击路径 & Indirect injection 可劫持执行链 & 环境输入与工具执行都需防护 \\
评测体系 & 传统 benchmark 集成成本高、标准弱 & 推动统一接口与可复现对抗评测 \\
防御策略 & 无银弹，必须 Defense-in-Depth & 多层拦截、动态策略、运行时约束 \\
治理方向 & 双用途风险要求持续监测 & 建立 observatory 与社区协作机制 \\
\bottomrule
\end{tabular}
\caption{Dawn Song 本讲可执行结论总览}
\end{table}

\begin{importantbox}{一句话结论}
Agentic AI 的安全问题不是 ``模型能不能拒答''，而是 ``系统在对抗环境下能否持续做对事，并阻止高危错误动作发生''。
\end{importantbox}

\subsection{延伸阅读}
\begin{itemize}
    \item Dawn Song 团队关于 Agent 安全与自动化红队的近期论文与项目主页。
    \item LLM/Agent 的 prompt injection、indirect injection 与 tool abuse 研究综述。
    \item 运行时策略执行与 least privilege 在智能体系统中的实践报告。
    \item Agent evaluation 标准化方向（MCP 风格接口、统一评测协议、arena 评测）。
    \item Frontier AI 与网络安全双用途风险治理相关报告与社区框架。
\end{itemize}

\end{document}
